\usepackage{fontspec}
\usepackage{polyglossia}

\setdefaultlanguage{arabic}
\setotherlanguage{english}

% Amiri remains the Arabic script font. LuaHBTeX/HarfBuzz preserves the
% logical Unicode text while shaping Arabic glyphs for display. Disable only
% common/contextual/discretionary ligatures so incidental Latin words in Arabic
% paragraphs do not reorder fi/fl/iff sequences. Required Arabic shaping
% features remain enabled by Script=Arabic.
\newfontfamily\arabicfont[
  Script=Arabic,
  Renderer=HarfBuzz,
  Scale=1.12,
  Ligatures=NoCommon,
  RawFeature={-liga;-clig;-dlig}
]{Amiri}

% Keep explicitly English text on an independent Latin font with the same
% defensive ligature policy.
\newfontfamily\englishfont[
  Script=Latin,
  Renderer=HarfBuzz,
  Language=English,
  Ligatures=NoCommon,
  RawFeature={-liga;-clig;-dlig}
]{Latin Modern Roman}

\newfontfamily\arabicfonttt[
  Script=Arabic,
  Renderer=HarfBuzz,
  Scale=1.05,
  Ligatures=NoCommon,
  RawFeature={-liga;-clig;-dlig}
]{Amiri}

\usepackage{amsmath,amssymb,amsthm,mathtools}
\usepackage{geometry}
\usepackage{xcolor}
\usepackage{xurl}
\usepackage{hyperref}
\usepackage{bookmark}
\usepackage{enumitem}
\usepackage{microtype}
\usepackage{csquotes}
\DeclareQuoteAlias[american]{english}{arabic}
\usepackage{fancyhdr}

% The named indexes are built by the repository's UTF-8-safe Python index
% generator between LaTeX passes. Keep imakeidx from invoking legacy makeindex.
\usepackage[noautomatic]{imakeidx}

% فهارس الموسوعة المرجعية.
\makeindex[
  name=people,
  title={فهرس العلماء},
  columns=2,
  intoc
]
\makeindex[
  name=theorems,
  title={فهرس النظريات والنتائج},
  columns=2,
  intoc
]
\makeindex[
  name=symbols,
  title={فهرس الرموز والمصطلحات},
  columns=2,
  intoc
]

% واجهة موحدة تمنع خلط أنواع الفهارس داخل الفصول.
\newcommand{\personindex}[1]{\index[people]{#1}}
\newcommand{\theoremindex}[1]{\index[theorems]{#1}}
\newcommand{\symbolindex}[1]{\index[symbols]{#1}}
% Polyglossia's \textenglish is a text-mode command. Preserve its normal
% behavior while allowing guarded use inside displayed mathematics.
\let\polyglossiatextenglish\textenglish
\renewcommand{\textenglish}[1]{%
  \ifmmode
    \text{\polyglossiatextenglish{#1}}%
  \else
    \polyglossiatextenglish{#1}%
  \fi
}

\geometry{a4paper,margin=27mm}

\usepackage[
  backend=biber,
  style=alphabetic,
  sorting=nyt,
  language=english
]{biblatex}

% LTR direction alone does not select a Latin font. Force every printed
% bibliography entry through the same ligature-safe English font.
\renewcommand*{\bibfont}{\footnotesize\englishfont}

% ---------------------------------------------------------------------------
% Visual layer: every figure in this encyclopedia is either a plot of a
% closed-form/finite arithmetic object computed by scripts/build_figure_data.py
% (imported below), or a schematic diagram of a definition already stated in
% the surrounding prose. No figure introduces a claim the text does not
% already make. Tick and axis text render through the Latin font used for
% every other piece of embedded English/math typography in the book.
\usepackage{tikz}
\usepackage{pgfplots}
\pgfplotsset{compat=1.18}
\usetikzlibrary{arrows.meta,positioning,decorations.pathreplacing,calc}
\pgfplotsset{
  every axis/.append style={
    font=\englishfont,
    label style={font=\englishfont},
    tick label style={font=\englishfont},
    legend style={font=\englishfont},
  }
}
\newenvironment{antfigure}
  {\begin{figure}[htbp]\centering}
  {\end{figure}}
% !TeX root = main.tex
% Generated by scripts/build_figure_data.py. Do not edit by hand.
% Each macro expands to a complete \addplot command; its single argument is
% the pgfplots option list applied at the call site. Expanding a whole command
% keeps the coordinate list literal at the moment pgfplots parses it.
\newcommand{\antPlotPsiExact}[1]{\addplot[#1] coordinates {%
(2.000000,-2.000000) (2.000000,-1.306853) (3.000000,-2.306853) (3.000000,-1.208241) (4.000000,-2.208241) (4.000000,-1.515093)%
(5.000000,-2.515093) (5.000000,-0.905655) (7.000000,-2.905655) (7.000000,-0.959745) (8.000000,-1.959745) (8.000000,-1.266598)%
(9.000000,-2.266598) (9.000000,-1.167986) (11.000000,-3.167986) (11.000000,-0.770091) (13.000000,-2.770091) (13.000000,-0.205141)%
(16.000000,-3.205141) (16.000000,-2.511994) (17.000000,-3.511994) (17.000000,-0.678781) (19.000000,-2.678781) (19.000000,0.265658)%
(23.000000,-3.734342) (23.000000,-0.598847) (25.000000,-2.598847) (25.000000,-0.989410) (27.000000,-2.989410) (27.000000,-1.890797)%
(29.000000,-3.890797) (29.000000,-0.523501) (31.000000,-2.523501) (31.000000,0.910486) (32.000000,-0.089514) (32.000000,0.603633)%
(37.000000,-4.396367) (37.000000,-0.785449) (41.000000,-4.785449) (41.000000,-1.071877) (43.000000,-3.071877) (43.000000,0.689323)%
(47.000000,-3.310677) (47.000000,0.539471) (49.000000,-1.460529) (49.000000,0.485381) (53.000000,-3.514619) (53.000000,0.455673)%
(59.000000,-5.544327) (59.000000,-1.466790) (61.000000,-3.466790) (61.000000,0.644084) (64.000000,-2.355916) (64.000000,-1.662769)%
(67.000000,-4.662769) (67.000000,-0.458076) (71.000000,-4.458076) (71.000000,-0.195396) (73.000000,-2.195396) (73.000000,2.095063)%
(79.000000,-3.904937) (79.000000,0.464511) (81.000000,-1.535489) (81.000000,-0.436877) (83.000000,-2.436877) (83.000000,1.981964)%
(89.000000,-4.018036) (89.000000,0.470600) (97.000000,-7.529400) (97.000000,-2.954689) (100.000000,-5.954689)
};}
\newcommand{\antPlotPsiExplicit}[1]{\addplot[#1] coordinates {%
(2.000000,-1.678421) (2.233333,-1.612761) (2.466667,-1.809247) (2.700000,-1.934395) (2.933333,-2.081179) (3.166667,-1.299347)%
(3.400000,-1.627169) (3.633333,-1.856406) (3.866667,-2.020743) (4.100000,-1.752100) (4.333333,-1.754443) (4.566667,-2.163539)%
(4.800000,-2.301157) (5.033333,-1.598185) (5.266667,-1.069037) (5.500000,-1.337598) (5.733333,-1.760054) (5.966667,-1.844992)%
(6.200000,-1.981649) (6.433333,-2.463310) (6.666667,-2.713088) (6.900000,-2.222361) (7.133333,-1.484704) (7.366667,-1.255439)%
(7.600000,-1.506250) (7.833333,-1.698107) (8.066667,-1.643162) (8.300000,-1.601098) (8.533333,-1.737788) (8.766667,-1.884099)%
(9.000000,-1.827337) (9.233333,-1.614973) (9.466667,-1.497534) (9.700000,-1.662681) (9.933333,-2.076427) (10.166667,-2.530738)%
(10.400000,-2.788435) (10.633333,-2.705976) (10.866667,-2.295696) (11.100000,-1.726834) (11.333333,-1.258936) (11.566667,-1.116665)%
(11.800000,-1.362195) (12.033333,-1.846933) (12.266667,-2.285687) (12.500000,-2.415486) (12.733333,-2.141983) (12.966667,-1.584915)%
(13.200000,-1.001470) (13.433333,-0.642074) (13.666667,-0.626889) (13.900000,-0.907953) (14.133333,-1.324337) (14.366667,-1.705286)%
(14.600000,-1.957994) (14.833333,-2.096103) (15.066667,-2.204369) (15.300000,-2.369474) (15.533333,-2.619564) (15.766667,-2.903763)%
(16.000000,-3.117808) (16.233333,-3.158089) (16.466667,-2.974779) (16.700000,-2.598674) (16.933333,-2.131237) (17.166667,-1.704703)%
(17.400000,-1.431071) (17.633333,-1.361360) (17.866667,-1.470114) (18.100000,-1.668861) (18.333333,-1.841127) (18.566667,-1.884735)%
(18.800000,-1.746282) (19.033333,-1.437134) (19.266667,-1.027622) (19.500000,-0.623476) (19.733333,-0.333512) (19.966667,-0.239007)%
(20.200000,-0.373338) (20.433333,-0.716329) (20.666667,-1.203035) (20.900000,-1.742802) (21.133333,-2.242348) (21.366667,-2.626500)%
(21.600000,-2.851848) (21.833333,-2.911097) (22.066667,-2.828591) (22.300000,-2.649523) (22.533333,-2.426460) (22.766667,-2.206787)%
(23.000000,-2.023830) (23.233333,-1.893029) (23.466667,-1.813106) (23.700000,-1.771006) (23.933333,-1.748713) (24.166667,-1.729942)%
(24.400000,-1.705104) (24.633333,-1.673583) (24.866667,-1.643184) (25.100000,-1.627275) (25.333333,-1.640616) (25.566667,-1.695038)%
(25.800000,-1.795993) (26.033333,-1.940718) (26.266667,-2.118285) (26.500000,-2.311388) (26.733333,-2.499361) (26.966667,-2.661733)%
(27.200000,-2.781566) (27.433333,-2.847955) (27.666667,-2.857271) (27.900000,-2.813029) (28.133333,-2.724508) (28.366667,-2.604480)%
(28.600000,-2.466538) (28.833333,-2.322536) (29.066667,-2.180594) (29.300000,-2.043993) (29.533333,-1.911092) (29.766667,-1.776228)%
(30.000000,-1.631393) (30.233333,-1.468369) (30.466667,-1.280922) (30.700000,-1.066718) (30.933333,-0.828614) (31.166667,-0.575154)%
(31.400000,-0.320175) (31.633333,-0.081583) (31.866667,0.120520) (32.100000,0.266122) (32.333333,0.337672) (32.566667,0.322263)%
(32.800000,0.213333) (33.033333,0.011690) (33.266667,-0.274245) (33.500000,-0.629007) (33.733333,-1.031386) (33.966667,-1.456304)%
(34.200000,-1.877060) (34.433333,-2.267703) (34.666667,-2.605287) (34.900000,-2.871801) (35.133333,-3.055578) (35.366667,-3.152093)%
(35.600000,-3.164077) (35.833333,-3.100990) (36.066667,-2.977914) (36.300000,-2.814004) (36.533333,-2.630656) (36.766667,-2.449569)%
(37.000000,-2.290876) (37.233333,-2.171503) (37.466667,-2.103882) (37.700000,-2.095108) (37.933333,-2.146581) (38.166667,-2.254144)%
(38.400000,-2.408658) (38.633333,-2.596960) (38.866667,-2.803086) (39.100000,-3.009657) (39.333333,-3.199306) (39.566667,-3.356023)%
(39.800000,-3.466332) (40.033333,-3.520210) (40.266667,-3.511704) (40.500000,-3.439206) (40.733333,-3.305405) (40.966667,-3.116923)%
(41.200000,-2.883694) (41.433333,-2.618146) (41.666667,-2.334249) (41.900000,-2.046514) (42.133333,-1.769016) (42.366667,-1.514504)%
(42.600000,-1.293648) (42.833333,-1.114479) (43.066667,-0.982035) (43.300000,-0.898223) (43.533333,-0.861899) (43.766667,-0.869138)%
(44.000000,-0.913663) (44.233333,-0.987399) (44.466667,-1.081105) (44.700000,-1.185032) (44.933333,-1.289575) (45.166667,-1.385864)%
(45.400000,-1.466275) (45.633333,-1.524822) (45.866667,-1.557424) (46.100000,-1.562033) (46.333333,-1.538629) (46.566667,-1.489087)%
(46.800000,-1.416940) (47.033333,-1.327055) (47.266667,-1.225248) (47.500000,-1.117865) (47.733333,-1.011362) (47.966667,-0.911906)%
(48.200000,-0.825010) (48.433333,-0.755242) (48.666667,-0.706002) (48.900000,-0.679382) (49.133333,-0.676117) (49.366667,-0.695613)%
(49.600000,-0.736061) (49.833333,-0.794609) (50.066667,-0.867588) (50.300000,-0.950779) (50.533333,-1.039690) (50.766667,-1.129846)%
(51.000000,-1.217057) (51.233333,-1.297657) (51.466667,-1.368712) (51.700000,-1.428167) (51.933333,-1.474940) (52.166667,-1.508965)%
(52.400000,-1.531166) (52.633333,-1.543382) (52.866667,-1.548244) (53.100000,-1.549005) (53.333333,-1.549346) (53.566667,-1.553160)%
(53.800000,-1.564329) (54.033333,-1.586495) (54.266667,-1.622864) (54.500000,-1.676011) (54.733333,-1.747735) (54.966667,-1.838941)%
(55.200000,-1.949572) (55.433333,-2.078580) (55.666667,-2.223949) (55.900000,-2.382752) (56.133333,-2.551259) (56.366667,-2.725073)%
(56.600000,-2.899296) (56.833333,-3.068715) (57.066667,-3.228006) (57.300000,-3.371933) (57.533333,-3.495554) (57.766667,-3.594410)%
(58.000000,-3.664693) (58.233333,-3.703397) (58.466667,-3.708433) (58.700000,-3.678711) (58.933333,-3.614190) (59.166667,-3.515889)%
(59.400000,-3.385859) (59.633333,-3.227124) (59.866667,-3.043588) (60.100000,-2.839914) (60.333333,-2.621380) (60.566667,-2.393720)%
(60.800000,-2.162943) (61.033333,-1.935158) (61.266667,-1.716388) (61.500000,-1.512399) (61.733333,-1.328528) (61.966667,-1.169542)%
(62.200000,-1.039498) (62.433333,-0.941646) (62.666667,-0.878341) (62.900000,-0.850993) (63.133333,-0.860038) (63.366667,-0.904945)%
(63.600000,-0.984243) (63.833333,-1.095574) (64.066667,-1.235775) (64.300000,-1.400974) (64.533333,-1.586704) (64.766667,-1.788030)%
(65.000000,-1.999689) (65.233333,-2.216226) (65.466667,-2.432141) (65.700000,-2.642024) (65.933333,-2.840690) (66.166667,-3.023303)%
(66.400000,-3.185482) (66.633333,-3.323402) (66.866667,-3.433865) (67.100000,-3.514366) (67.333333,-3.563131) (67.566667,-3.579138)%
(67.800000,-3.562121) (68.033333,-3.512554) (68.266667,-3.431615) (68.500000,-3.321144) (68.733333,-3.183571) (68.966667,-3.021849)%
(69.200000,-2.839368) (69.433333,-2.639867) (69.666667,-2.427333) (69.900000,-2.205910) (70.133333,-1.979801) (70.366667,-1.753173)%
(70.600000,-1.530068) (70.833333,-1.314321) (71.066667,-1.109485) (71.300000,-0.918768) (71.533333,-0.744975) (71.766667,-0.590466)%
(72.000000,-0.457127) (72.233333,-0.346346) (72.466667,-0.259010) (72.700000,-0.195503) (72.933333,-0.155722) (73.166667,-0.139102)%
(73.400000,-0.144647) (73.633333,-0.170968) (73.866667,-0.216333) (74.100000,-0.278716) (74.333333,-0.355853) (74.566667,-0.445296)%
(74.800000,-0.544477) (75.033333,-0.650759) (75.266667,-0.761499) (75.500000,-0.874093) (75.733333,-0.986034) (75.966667,-1.094952)%
(76.200000,-1.198657) (76.433333,-1.295169) (76.666667,-1.382754) (76.900000,-1.459934) (77.133333,-1.525516) (77.366667,-1.578587)%
(77.600000,-1.618530) (77.833333,-1.645008) (78.066667,-1.657967) (78.300000,-1.657613) (78.533333,-1.644399) (78.766667,-1.619002)%
(79.000000,-1.582296) (79.233333,-1.535330) (79.466667,-1.479290) (79.700000,-1.415479) (79.933333,-1.345278) (80.166667,-1.270123)%
(80.400000,-1.191468) (80.633333,-1.110766) (80.866667,-1.029437) (81.100000,-0.948844) (81.333333,-0.870277) (81.566667,-0.794931)%
(81.800000,-0.723893) (82.033333,-0.658128) (82.266667,-0.598473) (82.500000,-0.545630) (82.733333,-0.500164) (82.966667,-0.462504)%
(83.200000,-0.432945) (83.433333,-0.411655) (83.666667,-0.398683) (83.900000,-0.393967) (84.133333,-0.397348) (84.366667,-0.408581)%
(84.600000,-0.427346) (84.833333,-0.453265) (85.066667,-0.485918) (85.300000,-0.524849) (85.533333,-0.569587) (85.766667,-0.619656)%
(86.000000,-0.674585) (86.233333,-0.733921) (86.466667,-0.797237) (86.700000,-0.864142) (86.933333,-0.934282) (87.166667,-1.007353)%
(87.400000,-1.083095) (87.633333,-1.161300) (87.866667,-1.241809) (88.100000,-1.324514) (88.333333,-1.409350) (88.566667,-1.496293)%
(88.800000,-1.585358) (89.033333,-1.676590) (89.266667,-1.770055) (89.500000,-1.865838) (89.733333,-1.964032) (89.966667,-2.064730)%
(90.200000,-2.168017) (90.433333,-2.273965) (90.666667,-2.382621) (90.900000,-2.494004) (91.133333,-2.608096) (91.366667,-2.724839)%
(91.600000,-2.844127) (91.833333,-2.965805) (92.066667,-3.089664) (92.300000,-3.215441) (92.533333,-3.342818) (92.766667,-3.471420)%
(93.000000,-3.600819) (93.233333,-3.730535) (93.466667,-3.860040) (93.700000,-3.988763) (93.933333,-4.116091) (94.166667,-4.241380)%
(94.400000,-4.363958) (94.633333,-4.483133) (94.866667,-4.598200) (95.100000,-4.708449) (95.333333,-4.813171) (95.566667,-4.911668)%
(95.800000,-5.003258) (96.033333,-5.087287) (96.266667,-5.163129) (96.500000,-5.230201) (96.733333,-5.287961) (96.966667,-5.335922)%
(97.200000,-5.373649) (97.433333,-5.400770) (97.666667,-5.416976) (97.900000,-5.422024) (98.133333,-5.415740) (98.366667,-5.398022)%
(98.600000,-5.368834) (98.833333,-5.328216) (99.066667,-5.276271) (99.300000,-5.213175) (99.533333,-5.139165) (99.766667,-5.054540)%
(100.000000,-4.959659)
};}
\newcommand{\antPlotZeroFreeClassical}[1]{\addplot[#1] coordinates {%
(1.000000,0.565706) (1.181250,0.632343) (1.362500,0.681252) (1.543750,0.718676) (1.725000,0.748235) (1.906250,0.772173)%
(2.087500,0.791955) (2.268750,0.808575) (2.450000,0.822737) (2.631250,0.834947) (2.812500,0.845584) (2.993750,0.854933)%
(3.175000,0.863214) (3.356250,0.870601) (3.537500,0.877231) (3.718750,0.883215) (3.900000,0.888642) (4.081250,0.893588)%
(4.262500,0.898113) (4.443750,0.902268) (4.625000,0.906098) (4.806250,0.909640) (4.987500,0.912923) (5.168750,0.915977)%
(5.350000,0.918823) (5.531250,0.921483) (5.712500,0.923975) (5.893750,0.926313) (6.075000,0.928511) (6.256250,0.930582)%
(6.437500,0.932537) (6.618750,0.934384) (6.800000,0.936133) (6.981250,0.937791) (7.162500,0.939366) (7.343750,0.940862)%
(7.525000,0.942286) (7.706250,0.943644) (7.887500,0.944939) (8.068750,0.946176) (8.250000,0.947358) (8.431250,0.948490)%
(8.612500,0.949574) (8.793750,0.950613) (8.975000,0.951611) (9.156250,0.952569) (9.337500,0.953489) (9.518750,0.954375)%
(9.700000,0.955227) (9.881250,0.956049) (10.062500,0.956840) (10.243750,0.957604) (10.425000,0.958341) (10.606250,0.959053)%
(10.787500,0.959741) (10.968750,0.960406) (11.150000,0.961050) (11.331250,0.961673) (11.512500,0.962276) (11.693750,0.962861)%
(11.875000,0.963428) (12.056250,0.963978) (12.237500,0.964511) (12.418750,0.965029) (12.600000,0.965532) (12.781250,0.966021)%
(12.962500,0.966496) (13.143750,0.966958) (13.325000,0.967408) (13.506250,0.967845) (13.687500,0.968271) (13.868750,0.968685)%
(14.050000,0.969089) (14.231250,0.969483) (14.412500,0.969867) (14.593750,0.970241) (14.775000,0.970606) (14.956250,0.970962)%
(15.137500,0.971310) (15.318750,0.971649) (15.500000,0.971981) (15.681250,0.972305) (15.862500,0.972621) (16.043750,0.972931)%
(16.225000,0.973233) (16.406250,0.973529) (16.587500,0.973818) (16.768750,0.974101) (16.950000,0.974378) (17.131250,0.974649)%
(17.312500,0.974914) (17.493750,0.975174) (17.675000,0.975429) (17.856250,0.975678) (18.037500,0.975923) (18.218750,0.976162)%
(18.400000,0.976397) (18.581250,0.976627) (18.762500,0.976853) (18.943750,0.977075) (19.125000,0.977292) (19.306250,0.977505)%
(19.487500,0.977714) (19.668750,0.977920) (19.850000,0.978121) (20.031250,0.978319) (20.212500,0.978514) (20.393750,0.978705)%
(20.575000,0.978892) (20.756250,0.979076) (20.937500,0.979258) (21.118750,0.979436) (21.300000,0.979611) (21.481250,0.979783)%
(21.662500,0.979952) (21.843750,0.980118) (22.025000,0.980282) (22.206250,0.980443) (22.387500,0.980601) (22.568750,0.980757)%
(22.750000,0.980910) (22.931250,0.981061) (23.112500,0.981210) (23.293750,0.981356) (23.475000,0.981500) (23.656250,0.981641)%
(23.837500,0.981781) (24.018750,0.981919) (24.200000,0.982054) (24.381250,0.982187) (24.562500,0.982319) (24.743750,0.982448)%
(24.925000,0.982576) (25.106250,0.982702) (25.287500,0.982826) (25.468750,0.982948) (25.650000,0.983068) (25.831250,0.983187)%
(26.012500,0.983304) (26.193750,0.983420) (26.375000,0.983534) (26.556250,0.983646) (26.737500,0.983757) (26.918750,0.983866)%
(27.100000,0.983974) (27.281250,0.984081) (27.462500,0.984186) (27.643750,0.984290) (27.825000,0.984392) (28.006250,0.984493)%
(28.187500,0.984593) (28.368750,0.984691) (28.550000,0.984788) (28.731250,0.984884) (28.912500,0.984979) (29.093750,0.985073)%
(29.275000,0.985165) (29.456250,0.985256) (29.637500,0.985346) (29.818750,0.985436) (30.000000,0.985524)
};}
\newcommand{\antPlotZeroFreeKorobov}[1]{\addplot[#1] coordinates {%
(1.000000,0.390750) (1.181250,0.486894) (1.362500,0.553757) (1.543750,0.603354) (1.725000,0.641809) (1.906250,0.672611)%
(2.087500,0.697908) (2.268750,0.719100) (2.450000,0.737143) (2.631250,0.752711) (2.812500,0.766299) (2.993750,0.778273)%
(3.175000,0.788915) (3.356250,0.798442) (3.537500,0.807028) (3.718750,0.814809) (3.900000,0.821898) (4.081250,0.828387)%
(4.262500,0.834351) (4.443750,0.839854) (4.625000,0.844949) (4.806250,0.849682) (4.987500,0.854092) (5.168750,0.858211)%
(5.350000,0.862068) (5.531250,0.865690) (5.712500,0.869096) (5.893750,0.872308) (6.075000,0.875341) (6.256250,0.878212)%
(6.437500,0.880932) (6.618750,0.883514) (6.800000,0.885969) (6.981250,0.888306) (7.162500,0.890534) (7.343750,0.892661)%
(7.525000,0.894693) (7.706250,0.896638) (7.887500,0.898500) (8.068750,0.900286) (8.250000,0.901999) (8.431250,0.903645)%
(8.612500,0.905228) (8.793750,0.906751) (8.975000,0.908218) (9.156250,0.909631) (9.337500,0.910995) (9.518750,0.912311)%
(9.700000,0.913583) (9.881250,0.914812) (10.062500,0.916000) (10.243750,0.917151) (10.425000,0.918265) (10.606250,0.919344)%
(10.787500,0.920391) (10.968750,0.921406) (11.150000,0.922391) (11.331250,0.923347) (11.512500,0.924277) (11.693750,0.925180)%
(11.875000,0.926058) (12.056250,0.926912) (12.237500,0.927744) (12.418750,0.928553) (12.600000,0.929342) (12.781250,0.930110)%
(12.962500,0.930859) (13.143750,0.931589) (13.325000,0.932301) (13.506250,0.932996) (13.687500,0.933675) (13.868750,0.934337)%
(14.050000,0.934984) (14.231250,0.935617) (14.412500,0.936234) (14.593750,0.936839) (14.775000,0.937429) (14.956250,0.938007)%
(15.137500,0.938573) (15.318750,0.939126) (15.500000,0.939668) (15.681250,0.940198) (15.862500,0.940718) (16.043750,0.941227)%
(16.225000,0.941726) (16.406250,0.942215) (16.587500,0.942695) (16.768750,0.943165) (16.950000,0.943626) (17.131250,0.944078)%
(17.312500,0.944522) (17.493750,0.944958) (17.675000,0.945386) (17.856250,0.945805) (18.037500,0.946218) (18.218750,0.946623)%
(18.400000,0.947021) (18.581250,0.947412) (18.762500,0.947796) (18.943750,0.948173) (19.125000,0.948544) (19.306250,0.948909)%
(19.487500,0.949268) (19.668750,0.949621) (19.850000,0.949969) (20.031250,0.950310) (20.212500,0.950646) (20.393750,0.950977)%
(20.575000,0.951303) (20.756250,0.951623) (20.937500,0.951939) (21.118750,0.952250) (21.300000,0.952556) (21.481250,0.952857)%
(21.662500,0.953154) (21.843750,0.953447) (22.025000,0.953735) (22.206250,0.954019) (22.387500,0.954299) (22.568750,0.954575)%
(22.750000,0.954847) (22.931250,0.955115) (23.112500,0.955380) (23.293750,0.955640) (23.475000,0.955898) (23.656250,0.956151)%
(23.837500,0.956402) (24.018750,0.956648) (24.200000,0.956892) (24.381250,0.957133) (24.562500,0.957370) (24.743750,0.957604)%
(24.925000,0.957835) (25.106250,0.958063) (25.287500,0.958289) (25.468750,0.958511) (25.650000,0.958731) (25.831250,0.958947)%
(26.012500,0.959162) (26.193750,0.959373) (26.375000,0.959582) (26.556250,0.959789) (26.737500,0.959993) (26.918750,0.960194)%
(27.100000,0.960393) (27.281250,0.960590) (27.462500,0.960785) (27.643750,0.960977) (27.825000,0.961167) (28.006250,0.961355)%
(28.187500,0.961540) (28.368750,0.961724) (28.550000,0.961905) (28.731250,0.962085) (28.912500,0.962262) (29.093750,0.962438)%
(29.275000,0.962612) (29.456250,0.962783) (29.637500,0.962953) (29.818750,0.963121) (30.000000,0.963287)
};}
\newcommand{\antPlotPairCorrelation}[1]{\addplot[#1] coordinates {%
(0.000000,0.000000) (0.005000,0.000082) (0.010000,0.000329) (0.015000,0.000740) (0.020000,0.001315) (0.025000,0.002054)%
(0.030000,0.002957) (0.035000,0.004024) (0.040000,0.005253) (0.045000,0.006644) (0.050000,0.008198) (0.055000,0.009912)%
(0.060000,0.011788) (0.065000,0.013823) (0.070000,0.016017) (0.075000,0.018369) (0.080000,0.020879) (0.085000,0.023544)%
(0.090000,0.026366) (0.095000,0.029341) (0.100000,0.032469) (0.105000,0.035749) (0.110000,0.039179) (0.115000,0.042758)%
(0.120000,0.046485) (0.125000,0.050359) (0.130000,0.054377) (0.135000,0.058538) (0.140000,0.062841) (0.145000,0.067284)%
(0.150000,0.071865) (0.155000,0.076582) (0.160000,0.081434) (0.165000,0.086419) (0.170000,0.091534) (0.175000,0.096778)%
(0.180000,0.102149) (0.185000,0.107645) (0.190000,0.113264) (0.195000,0.119003) (0.200000,0.124860) (0.205000,0.130833)%
(0.210000,0.136920) (0.215000,0.143119) (0.220000,0.149427) (0.225000,0.155841) (0.230000,0.162360) (0.235000,0.168981)%
(0.240000,0.175702) (0.245000,0.182519) (0.250000,0.189431) (0.255000,0.196434) (0.260000,0.203527) (0.265000,0.210706)%
(0.270000,0.217969) (0.275000,0.225314) (0.280000,0.232737) (0.285000,0.240236) (0.290000,0.247808) (0.295000,0.255450)%
(0.300000,0.263160) (0.305000,0.270935) (0.310000,0.278772) (0.315000,0.286668) (0.320000,0.294621) (0.325000,0.302627)%
(0.330000,0.310684) (0.335000,0.318789) (0.340000,0.326939) (0.345000,0.335131) (0.350000,0.343362) (0.355000,0.351630)%
(0.360000,0.359932) (0.365000,0.368264) (0.370000,0.376624) (0.375000,0.385009) (0.380000,0.393417) (0.385000,0.401844)%
(0.390000,0.410288) (0.395000,0.418745) (0.400000,0.427213) (0.405000,0.435690) (0.410000,0.444172) (0.415000,0.452656)%
(0.420000,0.461141) (0.425000,0.469622) (0.430000,0.478098) (0.435000,0.486566) (0.440000,0.495023) (0.445000,0.503466)%
(0.450000,0.511893) (0.455000,0.520301) (0.460000,0.528688) (0.465000,0.537051) (0.470000,0.545388) (0.475000,0.553695)%
(0.480000,0.561972) (0.485000,0.570214) (0.490000,0.578421) (0.495000,0.586588) (0.500000,0.594715) (0.505000,0.602799)%
(0.510000,0.610837) (0.515000,0.618828) (0.520000,0.626769) (0.525000,0.634658) (0.530000,0.642493) (0.535000,0.650271)%
(0.540000,0.657992) (0.545000,0.665652) (0.550000,0.673251) (0.555000,0.680785) (0.560000,0.688254) (0.565000,0.695655)%
(0.570000,0.702987) (0.575000,0.710247) (0.580000,0.717435) (0.585000,0.724549) (0.590000,0.731586) (0.595000,0.738547)%
(0.600000,0.745428) (0.605000,0.752229) (0.610000,0.758949) (0.615000,0.765585) (0.620000,0.772137) (0.625000,0.778603)%
(0.630000,0.784983) (0.635000,0.791275) (0.640000,0.797478) (0.645000,0.803592) (0.650000,0.809614) (0.655000,0.815544)%
(0.660000,0.821382) (0.665000,0.827127) (0.670000,0.832777) (0.675000,0.838332) (0.680000,0.843791) (0.685000,0.849155)%
(0.690000,0.854421) (0.695000,0.859591) (0.700000,0.864662) (0.705000,0.869635) (0.710000,0.874510) (0.715000,0.879286)%
(0.720000,0.883963) (0.725000,0.888541) (0.730000,0.893019) (0.735000,0.897398) (0.740000,0.901677) (0.745000,0.905857)%
(0.750000,0.909937) (0.755000,0.913917) (0.760000,0.917798) (0.765000,0.921581) (0.770000,0.925264) (0.775000,0.928848)%
(0.780000,0.932334) (0.785000,0.935723) (0.790000,0.939013) (0.795000,0.942207) (0.800000,0.945304) (0.805000,0.948305)%
(0.810000,0.951210) (0.815000,0.954020) (0.820000,0.956737) (0.825000,0.959359) (0.830000,0.961889) (0.835000,0.964327)%
(0.840000,0.966673) (0.845000,0.968929) (0.850000,0.971096) (0.855000,0.973174) (0.860000,0.975165) (0.865000,0.977068)%
(0.870000,0.978886) (0.875000,0.980620) (0.880000,0.982269) (0.885000,0.983837) (0.890000,0.985323) (0.895000,0.986728)%
(0.900000,0.988055) (0.905000,0.989304) (0.910000,0.990476) (0.915000,0.991573) (0.920000,0.992596) (0.925000,0.993547)%
(0.930000,0.994425) (0.935000,0.995234) (0.940000,0.995974) (0.945000,0.996646) (0.950000,0.997253) (0.955000,0.997794)%
(0.960000,0.998273) (0.965000,0.998690) (0.970000,0.999046) (0.975000,0.999344) (0.980000,0.999584) (0.985000,0.999768)%
(0.990000,0.999898) (0.995000,0.999975) (1.000000,1.000000) (1.005000,0.999975) (1.010000,0.999902) (1.015000,0.999782)%
(1.020000,0.999616) (1.025000,0.999406) (1.030000,0.999154) (1.035000,0.998861) (1.040000,0.998528) (1.045000,0.998158)%
(1.050000,0.997751) (1.055000,0.997309) (1.060000,0.996834) (1.065000,0.996326) (1.070000,0.995789) (1.075000,0.995222)%
(1.080000,0.994628) (1.085000,0.994007) (1.090000,0.993362) (1.095000,0.992694) (1.100000,0.992004) (1.105000,0.991293)%
(1.110000,0.990564) (1.115000,0.989817) (1.120000,0.989054) (1.125000,0.988276) (1.130000,0.987485) (1.135000,0.986681)%
(1.140000,0.985866) (1.145000,0.985042) (1.150000,0.984209) (1.155000,0.983370) (1.160000,0.982524) (1.165000,0.981674)%
(1.170000,0.980821) (1.175000,0.979965) (1.180000,0.979108) (1.185000,0.978251) (1.190000,0.977395) (1.195000,0.976541)%
(1.200000,0.975691) (1.205000,0.974844) (1.210000,0.974003) (1.215000,0.973169) (1.220000,0.972341) (1.225000,0.971522)%
(1.230000,0.970711) (1.235000,0.969911) (1.240000,0.969121) (1.245000,0.968343) (1.250000,0.967577) (1.255000,0.966825)%
(1.260000,0.966086) (1.265000,0.965362) (1.270000,0.964654) (1.275000,0.963961) (1.280000,0.963285) (1.285000,0.962627)%
(1.290000,0.961986) (1.295000,0.961363) (1.300000,0.960760) (1.305000,0.960176) (1.310000,0.959612) (1.315000,0.959068)%
(1.320000,0.958545) (1.325000,0.958043) (1.330000,0.957563) (1.335000,0.957105) (1.340000,0.956669) (1.345000,0.956255)%
(1.350000,0.955864) (1.355000,0.955496) (1.360000,0.955151) (1.365000,0.954829) (1.370000,0.954531) (1.375000,0.954257)%
(1.380000,0.954006) (1.385000,0.953779) (1.390000,0.953576) (1.395000,0.953397) (1.400000,0.953242) (1.405000,0.953111)%
(1.410000,0.953003) (1.415000,0.952919) (1.420000,0.952859) (1.425000,0.952823) (1.430000,0.952810) (1.435000,0.952820)%
(1.440000,0.952853) (1.445000,0.952909) (1.450000,0.952989) (1.455000,0.953090) (1.460000,0.953214) (1.465000,0.953359)%
(1.470000,0.953527) (1.475000,0.953716) (1.480000,0.953925) (1.485000,0.954156) (1.490000,0.954407) (1.495000,0.954678)%
(1.500000,0.954968) (1.505000,0.955278) (1.510000,0.955607) (1.515000,0.955954) (1.520000,0.956319) (1.525000,0.956701)%
(1.530000,0.957100) (1.535000,0.957516) (1.540000,0.957948) (1.545000,0.958396) (1.550000,0.958859) (1.555000,0.959336)%
(1.560000,0.959828) (1.565000,0.960333) (1.570000,0.960850) (1.575000,0.961381) (1.580000,0.961923) (1.585000,0.962477)%
(1.590000,0.963042) (1.595000,0.963616) (1.600000,0.964201) (1.605000,0.964794) (1.610000,0.965397) (1.615000,0.966007)%
(1.620000,0.966625) (1.625000,0.967249) (1.630000,0.967880) (1.635000,0.968516) (1.640000,0.969158) (1.645000,0.969804)%
(1.650000,0.970454) (1.655000,0.971108) (1.660000,0.971764) (1.665000,0.972423) (1.670000,0.973084) (1.675000,0.973746)%
(1.680000,0.974408) (1.685000,0.975071) (1.690000,0.975733) (1.695000,0.976394) (1.700000,0.977053) (1.705000,0.977711)%
(1.710000,0.978366) (1.715000,0.979018) (1.720000,0.979667) (1.725000,0.980311) (1.730000,0.980952) (1.735000,0.981587)%
(1.740000,0.982216) (1.745000,0.982840) (1.750000,0.983458) (1.755000,0.984069) (1.760000,0.984672) (1.765000,0.985268)%
(1.770000,0.985856) (1.775000,0.986436) (1.780000,0.987007) (1.785000,0.987569) (1.790000,0.988121) (1.795000,0.988663)%
(1.800000,0.989196) (1.805000,0.989718) (1.810000,0.990229) (1.815000,0.990729) (1.820000,0.991218) (1.825000,0.991695)%
(1.830000,0.992160) (1.835000,0.992613) (1.840000,0.993054) (1.845000,0.993483) (1.850000,0.993898) (1.855000,0.994301)%
(1.860000,0.994691) (1.865000,0.995067) (1.870000,0.995430) (1.875000,0.995779) (1.880000,0.996115) (1.885000,0.996437)%
(1.890000,0.996745) (1.895000,0.997040) (1.900000,0.997320) (1.905000,0.997586) (1.910000,0.997838) (1.915000,0.998076)%
(1.920000,0.998300) (1.925000,0.998510) (1.930000,0.998706) (1.935000,0.998887) (1.940000,0.999055) (1.945000,0.999208)%
(1.950000,0.999348) (1.955000,0.999474) (1.960000,0.999586) (1.965000,0.999684) (1.970000,0.999769) (1.975000,0.999840)%
(1.980000,0.999898) (1.985000,0.999943) (1.990000,0.999975) (1.995000,0.999994) (2.000000,1.000000) (2.005000,0.999994)%
(2.010000,0.999975) (2.015000,0.999945) (2.020000,0.999902) (2.025000,0.999848) (2.030000,0.999782) (2.035000,0.999705)%
(2.040000,0.999618) (2.045000,0.999519) (2.050000,0.999410) (2.055000,0.999291) (2.060000,0.999162) (2.065000,0.999023)%
(2.070000,0.998875) (2.075000,0.998718) (2.080000,0.998552) (2.085000,0.998377) (2.090000,0.998195) (2.095000,0.998004)%
(2.100000,0.997806) (2.105000,0.997601) (2.110000,0.997389) (2.115000,0.997170) (2.120000,0.996945) (2.125000,0.996714)%
(2.130000,0.996478) (2.135000,0.996236) (2.140000,0.995989) (2.145000,0.995738) (2.150000,0.995482) (2.155000,0.995223)%
(2.160000,0.994960) (2.165000,0.994694) (2.170000,0.994424) (2.175000,0.994153) (2.180000,0.993879) (2.185000,0.993603)%
(2.190000,0.993326) (2.195000,0.993047) (2.200000,0.992767) (2.205000,0.992487) (2.210000,0.992207) (2.215000,0.991927)%
(2.220000,0.991647) (2.225000,0.991368) (2.230000,0.991089) (2.235000,0.990813) (2.240000,0.990537) (2.245000,0.990264)%
(2.250000,0.989993) (2.255000,0.989724) (2.260000,0.989459) (2.265000,0.989196) (2.270000,0.988936) (2.275000,0.988680)%
(2.280000,0.988428) (2.285000,0.988181) (2.290000,0.987937) (2.295000,0.987698) (2.300000,0.987464) (2.305000,0.987235)%
(2.310000,0.987011) (2.315000,0.986793) (2.320000,0.986580) (2.325000,0.986373) (2.330000,0.986173) (2.335000,0.985978)%
(2.340000,0.985790) (2.345000,0.985609) (2.350000,0.985434) (2.355000,0.985267) (2.360000,0.985106) (2.365000,0.984953)%
(2.370000,0.984807) (2.375000,0.984668) (2.380000,0.984537) (2.385000,0.984413) (2.390000,0.984297) (2.395000,0.984189)%
(2.400000,0.984089) (2.405000,0.983997) (2.410000,0.983913) (2.415000,0.983837) (2.420000,0.983769) (2.425000,0.983709)%
(2.430000,0.983658) (2.435000,0.983614) (2.440000,0.983579) (2.445000,0.983552) (2.450000,0.983533) (2.455000,0.983523)%
(2.460000,0.983520) (2.465000,0.983526) (2.470000,0.983540) (2.475000,0.983561) (2.480000,0.983591) (2.485000,0.983629)%
(2.490000,0.983674) (2.495000,0.983728) (2.500000,0.983789) (2.505000,0.983857) (2.510000,0.983933) (2.515000,0.984017)%
(2.520000,0.984108) (2.525000,0.984206) (2.530000,0.984311) (2.535000,0.984423) (2.540000,0.984542) (2.545000,0.984667)%
(2.550000,0.984799) (2.555000,0.984938) (2.560000,0.985082) (2.565000,0.985233) (2.570000,0.985390) (2.575000,0.985552)%
(2.580000,0.985720) (2.585000,0.985893) (2.590000,0.986071) (2.595000,0.986255) (2.600000,0.986443) (2.605000,0.986636)%
(2.610000,0.986833) (2.615000,0.987034) (2.620000,0.987240) (2.625000,0.987449) (2.630000,0.987662) (2.635000,0.987878)%
(2.640000,0.988098) (2.645000,0.988320) (2.650000,0.988546) (2.655000,0.988773) (2.660000,0.989004) (2.665000,0.989236)%
(2.670000,0.989470) (2.675000,0.989706) (2.680000,0.989943) (2.685000,0.990182) (2.690000,0.990422) (2.695000,0.990662)%
(2.700000,0.990903) (2.705000,0.991145) (2.710000,0.991386) (2.715000,0.991628) (2.720000,0.991869) (2.725000,0.992110)%
(2.730000,0.992351) (2.735000,0.992590) (2.740000,0.992828) (2.745000,0.993065) (2.750000,0.993301) (2.755000,0.993535)%
(2.760000,0.993767) (2.765000,0.993997) (2.770000,0.994225) (2.775000,0.994450) (2.780000,0.994673) (2.785000,0.994893)%
(2.790000,0.995110) (2.795000,0.995324) (2.800000,0.995535) (2.805000,0.995742) (2.810000,0.995946) (2.815000,0.996146)%
(2.820000,0.996342) (2.825000,0.996534) (2.830000,0.996722) (2.835000,0.996905) (2.840000,0.997084) (2.845000,0.997259)%
(2.850000,0.997429) (2.855000,0.997594) (2.860000,0.997754) (2.865000,0.997910) (2.870000,0.998060) (2.875000,0.998205)%
(2.880000,0.998345) (2.885000,0.998479) (2.890000,0.998608) (2.895000,0.998732) (2.900000,0.998850) (2.905000,0.998962)%
(2.910000,0.999069) (2.915000,0.999170) (2.920000,0.999265) (2.925000,0.999355) (2.930000,0.999438) (2.935000,0.999516)%
(2.940000,0.999588) (2.945000,0.999655) (2.950000,0.999715) (2.955000,0.999770) (2.960000,0.999818) (2.965000,0.999861)%
(2.970000,0.999898) (2.975000,0.999930) (2.980000,0.999955) (2.985000,0.999975) (2.990000,0.999989) (2.995000,0.999997)%
(3.000000,1.000000)
};}
\newcommand{\antDataFareyOrderFive}{0.0000/0/1,72.0000/1/5,90.0000/1/4,120.0000/1/3,144.0000/2/5,180.0000/1/2,216.0000/3/5,240.0000/2/3,270.0000/3/4,288.0000/4/5,360.0000/1/1}


\addbibresource{manuscript/bibliography.bib}
\addbibresource{manuscript/chapter-13-bibliography.bib}
\addbibresource{manuscript/chapter-15-bibliography.bib}
\addbibresource{manuscript/chapter-16-bibliography.bib}
\addbibresource{manuscript/chapter-17-bibliography.bib}
\addbibresource{manuscript/chapter-18-bibliography.bib}
\addbibresource{manuscript/chapter-19-bibliography.bib}
\addbibresource{manuscript/chapter-20-bibliography.bib}
\addbibresource{manuscript/chapter-21-bibliography.bib}
\addbibresource{manuscript/chapter-22-bibliography.bib}
\addbibresource{manuscript/chapter-23-bibliography.bib}
\addbibresource{manuscript/chapter-24-bibliography.bib}
\addbibresource{manuscript/chapter-25-bibliography.bib}

\newtheorem{theorem}{مبرهنة}[chapter]
\newtheorem{lemma}[theorem]{لمّة}
\newtheorem{proposition}[theorem]{قضية}
\newtheorem{corollary}[theorem]{نتيجة}

\theoremstyle{definition}
\newtheorem{definition}[theorem]{تعريف}
\newtheorem{example}[theorem]{مثال}
\newtheorem{exercise}[theorem]{تمرين}
\newtheorem{conjecture}[theorem]{حدسية}
\newtheorem{openproblem}[theorem]{مسألة مفتوحة}

\theoremstyle{remark}
\newtheorem{remark}[theorem]{ملاحظة}

\newcommand{\N}{\mathbb{N}}
\newcommand{\Z}{\mathbb{Z}}
\newcommand{\Q}{\mathbb{Q}}
\newcommand{\R}{\mathbb{R}}
\newcommand{\C}{\mathbb{C}}
\newcommand{\RePart}{\operatorname{Re}}
\newcommand{\ImPart}{\operatorname{Im}}
\newcommand{\Res}{\operatorname{Res}}
\newcommand{\Li}{\operatorname{Li}}
\newcommand{\e}{\mathrm{e}}
\newcommand{\ii}{\mathrm{i}}

\newcommand{\resultbadge}[2]{%
  \noindent\colorbox{#1}{\strut\textbf{#2}}%
}
\newcommand{\resultid}[1]{%
  \noindent\texttt{#1}\par
}
\newcommand{\provedhere}{%
  \resultbadge{green!20}{مبرهن هنا}\par
}
\newcommand{\citedresult}[1]{%
  \resultbadge{blue!15}{مستشهد بها}\quad منقولة من المرجع #1.\par
}
\newcommand{\deferredresult}[1]{%
  \resultbadge{yellow!25}{برهان مؤجل}\quad #1\par
}
\newcommand{\conditionalresult}[1]{%
  \resultbadge{orange!25}{نتيجة مشروطة}\quad #1\par
}
\newcommand{\openresult}{%
  \resultbadge{red!20}{مسألة مفتوحة}\par
}

% بطاقة "أحدث ما وصل إليه الحدّ" للجبهات البحثية الحديثة (الفصول ٢٢--٢٦).
% تطبّق حرفيًا بروتوكول تحديث الخريطة في الفصل ٢٦: تاريخ قطع صريح، وحالة نشر،
% ومصدر أساسي، وتحديد الجملة/النتيجة المتأثرة. ليست بديلًا عن بحث محدّث مستقل؛
% إنما تجعل تاريخ آخر تحقق ظاهرًا للقارئ بدل عبارة "الأحدث" غير المؤرَّخة.
\newcommand{\frontierupdatebadge}{\resultbadge{cyan!20}{تحديث الجبهة}}
\newenvironment{frontiercard}[3]{%
  \par\medskip\noindent\frontierupdatebadge\par
  \noindent\textbf{تاريخ القطع:}~#1\hfill\textbf{حالة النشر:}~#2\par
  \noindent\textbf{المصدر الأساسي:}~#3\par
  \noindent\ignorespaces
}{%
  \par\medskip
}

\pagestyle{fancy}
\fancyhf{}
\fancyhead[LE,RO]{\thepage}
\fancyhead[RE]{\leftmark}
\fancyhead[LO]{\rightmark}
\setlength{\headheight}{15pt}

\hypersetup{
  colorlinks=true,
  linkcolor=blue!50!black,
  citecolor=green!40!black,
  urlcolor=blue!60!black,
  pdftitle={الموسوعة الشاملة في نظرية الأعداد التحليلية},
  pdfauthor={mesuef1974}
}
